% BHU PYQ -- Professional Exam Template
\documentclass[11pt,a4paper]{article}

\usepackage[a4paper,top=2.5cm,bottom=2.5cm,left=2.8cm,right=2.8cm,headheight=14pt]{geometry}
\usepackage{amsmath,amssymb}
\usepackage[T1]{fontenc}
\usepackage[utf8]{inputenc}
\usepackage{lmodern}
\usepackage{microtype}
\usepackage{fancyhdr}
\usepackage{enumitem}
\usepackage{listings}
\usepackage{hyperref}
\usepackage{lastpage}
\usepackage{parskip}

\hypersetup{colorlinks=true,urlcolor=black,linkcolor=black,
  pdftitle={BCS-301: Numerical Computing -- BHU PYQ}}

\pagestyle{fancy}
\fancyhf{}
\renewcommand{\headrulewidth}{0.4pt}
\renewcommand{\footrulewidth}{0.4pt}
\fancyhead[L]{\small   \textit{Banaras Hindu University}}
\fancyhead[R]{\small   \textit{BCS-301: Numerical Computing}}
\fancyfoot[C]{\small Page  \thepage\ of \pageref{LastPage}}


\newlist{parts}{enumerate}{2}
\setlist[parts,1]{label=(\alph*),leftmargin=*,itemsep=4pt,topsep=3pt}
\setlist[parts,2]{label=(\roman*),leftmargin=*,itemsep=2pt,topsep=2pt}

\newcommand{\pts}[1]{\hfill   \textit{[#1]}}


\begin{document}

\begin{center}
  {\large\bfseries BANARAS HINDU UNIVERSITY}\\[2pt]
  {
\normalsize B.Sc.\ (Hons.) Semester III Examination, 2016--17}\\[6pt]
  \rule{0.8\linewidth}{0.5pt}\\[6pt]
  {\large\bfseries Paper: BCS-301 --- Numerical Computing}\\[4pt]
  {
\normalsize Time: 3 Hours \hfill Maximum Marks: 70}
\end{center}

\rule{\linewidth}{0.4pt}

\smallskip



\noindent   \textit{Instructions: Answer   \textbf{five} questions including Question~1 which is compulsory. Use of calculator is allowed. All undefined symbols have their usual meanings. Each question carries   \textbf{14 marks}.}

\bigskip




\noindent   \textbf{Question 1.} (Compulsory --- attempt all parts)\pts{$2 imes7=14$}

\begin{parts}
  \item Explain different types of errors with an example.
  \item If the first significant figure of a number is $k$ and the number is correct to $n$ significant figures, prove that the relative error is less than $\dfrac{1}{2k  \times 10^{n-1}}$.
  \item Prove that the Newton--Raphson process has second-order (quadratic) convergence.
  \item Given that the equation $x^3 - 7 = 69$ has a root between 5 and 8, apply the method of Regula-Falsi (False Position) to determine it.
  \item Define Finite Differences. Explain Forward, Backward, and Central differences with examples.
  \item Evaluate $\Delta an^{-1}x$.
  \item State the conditions under which the Gauss--Seidel iteration method converges.
\end{parts}

\bigskip




\noindent   \textbf{Question 2.}\pts{$6+8=14$}

\begin{parts}
  \item Find a real root of the equation $x^3 - 2x - 5 = 0$ using the Bisection Method, correct to three decimal places.
  \item From the following table, estimate the number of students who obtained marks between 40 and 45 using Newton's Forward Interpolation Formula:
  
\begin{center}
  
\begin{tabular}{|c|c|c|c|c|c|}
  \hline
  Marks & 30--40 & 40--50 & 50--60 & 60--70 & 70--80 \\
  \hline
  No.\ of Students & 31 & 42 & 51 & 35 & 31 \\
  \hline
  \end{tabular}
  \end{center}
\end{parts}

\bigskip




\noindent   \textbf{Question 3.}\pts{$6+8=14$}

\begin{parts}
  \item Apply the Gauss Elimination method to solve the following system of equations, correct to three decimal places:
  \[
    x + y + z = 9, \quad 2x - 3y + 4z = 13, \quad 3x + 4y + 5z = 40
  \]
  \item Apply the Gauss--Jordan method to solve the following system, correct to three decimal places:
  \[
    x + 4y - z = -5, \quad x + y - 6z = -12, \quad 3x - y - z = 4
  \]
\end{parts}

\bigskip




\noindent   \textbf{Question 4.} The table gives the distance $y$ (in nautical miles) of the visible horizon for given heights $x$ (in feet) above the earth's surface:\pts{$6+8=14$}

\begin{center}

\begin{tabular}{|c|c|c|c|c|c|c|c|}
\hline
$x$ (ft)  & 100   & 150   & 200   & 250   & 300   & 350   & 400   \\
\hline
$y$ (mi.) & 10.63 & 13.03 & 15.04 & 16.81 & 18.42 & 19.90 & 21.27 \\
\hline
\end{tabular}
\end{center}
Find the values of $y$ when:

\begin{parts}
  \item $x = 218$ ft
  \item $x = 410$ ft
\end{parts}

\bigskip




\noindent   \textbf{Question 5.} Evaluate $\displaystyle\int_0^1 \frac{1}{1+x^2}\,dx$ using the following rules and compare with the actual value:\pts{$3+4+4+3=14$}

\begin{parts}
  \item Trapezoidal Rule
  \item Simpson's $ frac{1}{3}$ Rule
  \item Simpson's $ frac{3}{8}$ Rule
  \item Weddle's Rule
\end{parts}

\bigskip




\noindent   \textbf{Question 6.} Using the Runge--Kutta method of fourth order, solve $\dfrac{dy}{dx} = x + y$ with $y(0)=1$ at $x = 0.2$ and $x = 0.4$.\pts{14}

\bigskip




\noindent   \textbf{Question 7.}\pts{$6+8=14$}

\begin{parts}
  \item Apply Milne's method to find a solution of the differential equation $y' = x - y^2$ in the range $0 \le x \le 1$ for the boundary condition $y = 0$ at $x = 0$.
  \item Given $\dfrac{dx}{dy} = x^2(1+y)$ and $y(1)=1$,\ $y(1.1)=1.233$,\ $y(1.2)=1.548$,\ $y(1.3)=1.979$, evaluate $y(1.4)$ by the Adams--Bashforth method.
\end{parts}

\vfill
\rule{\linewidth}{0.4pt}

\begin{center}\small
\href{https://bhu-exam.vercel.app/}{Click here to visit our website}\\[4pt]\href{https://me-aryan.vercel.app/}{Click here to visit our contributor}\\[4pt]\href{https://quashan-me.vercel.app/}{Click here to visit our contributor}
\end{center}
\end{document}